% ==============================================================================
% Chapter 6: Conclusion & Future Outlook
% Target Page Budget: ~2 pages
% ==============================================================================
\chapter{Conclusion \& Future Outlook}
\label{chap:conclusion}

\section{Synthesis of the Preparatory Investigation}
\label{sec:conclusion_synthesis}
% [Drafting Guidance: ~1.0 page]
% - Reiterate the societal urgency of addressing social comparison and influencer exposure.
% - Revisit how the multi-agent architecture of SocialMirror-AI addresses fundamental theoretical gaps.
% - Summary of the theoretical and technical groundwork established in this preparatory specification.

\section{Anticipated Impact on Academic and Applied Domains}
\label{sec:conclusion_impact}
% [Drafting Guidance: ~0.5 pages]
% - Contributions to computational social science, affective computing, and agentic AI.
% - Potential translations into educational interventions, policy insights, and user empowerment tools.

\section{Concluding Thoughts}
\label{sec:conclusion_remarks}
% [Drafting Guidance: ~0.5 pages]
% - Final inspiring synthesis connecting algorithmic responsibility with user psychological agency.
The preparatory work detailed across this report demonstrates that autonomous multi-agent systems offer a transformative paradigm for understanding and mitigating the subtle cognitive impacts of algorithmic social media. By combining empirical psychological rigor with agentic cognitive modeling, \projectname{} aims to serve not merely as a passive analytical tool, but as an empowering reflective companion in the digital age.
